% Part: turing-machines
% Chapter: undecidability
% Section: universal-tm

\documentclass[../../../include/open-logic-section]{subfiles}

\begin{document}

\olfileid[hi]{tur}{und}{uni}
\olsection{सार्वत्रिक ट्यूरिंग मशीनें}

\olref[enu]{sec} में हमने चर्चा की थी कि प्रत्येक ट्यूरिंग मशीन का वर्णन
पूर्णांकों के एक परिमित अनुक्रम द्वारा कैसे किया जा सकता है। यह अनुक्रम
ट्यूरिंग मशीन की अवस्थाओं, वर्णमाला, आरंभिक अवस्था और निर्देशों को
कूटबद्ध करता है। हमने यह भी बताया था कि इन सभी वर्णनों का समुच्चय
!!{enumerable} है। चूँकि ऐसे वर्णनों का समुच्चय !!{denumerable} है,
इसका अर्थ है कि~$\Nat$ से इन वर्णनों तक !!a{surjective} फलन है। ऐसा
!!a{surjective} फलन, उदाहरणतः, कांटोर की टेढ़ी-मेढ़ी विधि से पाया जा
सकता है। इससे ट्यूरिंग मशीनों के सभी वर्णनों का परिगणन करने की रीति
मिलती है। ऐसा कोई परिगणन निश्चित कर लेने पर $1$वीं, $2$वीं, \dots,
$e$वीं ट्यूरिंग मशीन की बात करना सार्थक है। इन संख्याओं को
\emph{सूचकांक} कहते हैं।

\begin{defn}
यदि हमारे निश्चित परिगणन में $M$, $e$वीं ट्यूरिंग मशीन है, तो हम कहते
हैं कि $e$,~$M$ का एक \emph{सूचकांक} है। $M_e$ से हम $e$वीं ट्यूरिंग
मशीन निरूपित करते हैं।
\end{defn}

किसी मशीन के एक से अधिक सूचकांक हो सकते हैं। उदाहरणतः~$M$ के दो वर्णन
उसके निर्देशों को सूचीबद्ध करने के क्रम में भिन्न हो सकते हैं और इन
भिन्न वर्णनों के सूचकांक भी भिन्न होंगे।

महत्त्वपूर्ण बात यह है कि ट्यूरिंग मशीन के वर्णनों का परिगणन ऐसी रीति
से दिया जा सकता है कि मशीन~$M$ के सूचकांक से उसका वर्णन प्रभावी रूप से
संगणित किया जा सके और मशीन~$M$ के वर्णन से उसका कोई सूचकांक प्रभावी रूप
से संगणित किया जा सके। तब चर्च--ट्यूरिंग अभिधारणा के अनुसार ऐसी ट्यूरिंग
मशीन पाई जा सकती है जो सूचकांक~$e$ वाली ट्यूरिंग मशीन का वर्णन प्राप्त
करके उस अनुरूप वर्णन को अपनी पट्टी पर निर्गत के रूप में लिखती है। यह
वर्णन~$\TMstroke$ के खंडों का अनुक्रम होगा (जो~$M_e$ का वर्णन करने वाले
अनुक्रम के धन पूर्णांकों को निरूपित करता है)।

इसे देखते हुए अब यह पूछना स्वाभाविक है: ट्यूरिंग मशीन के सूचकांकों के
कौन-से फलन स्वयं ट्यूरिंग मशीनों द्वारा संगणनीय हैं? ट्यूरिंग मशीन के
सूचकांकों के किन गुणों का निर्णय ट्यूरिंग मशीनें कर सकती हैं? एक उदाहरण
लें: सूचकांक~$e$ को सूचकांक~$e$ वाली ट्यूरिंग मशीन की अवस्था-संख्या पर
भेजने वाला फलन किसी ट्यूरिंग मशीन द्वारा संगणनीय है। ऐसी ट्यूरिंग मशीन
यह करेगी: $e$ संख्या के~$\TMstroke$ का एक खंड रखने वाली पट्टी पर आरंभ
होकर वह पहले $e$ को उसके वर्णन में विकूटित करेगी। अब पट्टी पर वह वर्णन
~$\TMstroke$ के खंडों के अनुक्रम से निरूपित है। चूँकि इस अनुक्रम का पहला
!!{element} अवस्थाओं की संख्या है, इसलिए अब केवल पहले $\TMstroke$-खंड
को छोड़कर सब कुछ मिटाना और फिर रुकना है।

निम्न परिणाम उल्लेखनीय है:

\begin{thm}\ollabel{thm:universal-tm} एक \emph{सार्वत्रिक ट्यूरिंग
  मशीन}~$U$ है जो आगत $\tuple{e,n}$ पर आरंभ किए जाने पर
  \begin{enumerate}
    \item तभी रुकती है जब $M_e$ आगत~$n$ पर रुकती है, और
    \item यदि $M_e$ निर्गत $m$ के साथ रुकती है, तो~$U$ भी ऐसा करती है।
  \end{enumerate}
  अतः~$U$, फलन $f\colon \Nat \times \Nat \pto \Nat$ की संगणना करती
  है, जहाँ $f(e,n) = m$ तब है जब $M_e$, आगत~$n$ पर आरंभ होकर निर्गत
  ~$m$ के साथ रुकती है, और अन्यथा अपरिभाषित है।
\end{thm}

\begin{proof}
  वास्तव में~$U$ का निर्माण करना लगभग असंभव है, क्योंकि वह अत्यंत जटिल
  मशीन है। किंतु हम उसके कार्य की रूपरेखा बता सकते हैं और फिर
  चर्च--ट्यूरिंग अभिधारणा का आह्वान कर सकते हैं। आरंभ के समय~$U$ की पट्टी
  पर $e$ संख्या के~$\TMstroke$ का एक खंड और उसके बाद $n$ संख्या के
  ~ $\TMstroke$ का एक खंड होता है। वह पहले सूचकांक~$e$ को आगत~$n$ की
  दाईं ओर “विकूटित” करती है। इससे संख्याओं की ऐसी सूची (अर्थात्
  $\TMstroke$ के ऐसे खंड जो~$\TMblank$ से पृथक हैं) बनती है जो
  मशीन~$M_e$ के निर्देशों का वर्णन करती है। फिर~$U$, $M_e$ की आरंभिक
  अवस्था की संख्या और संख्या~$1$, $M_e$ के वर्णन की दाईं ओर पट्टी पर
  लिखती है। (इनका भी
  एकाधारी रूप में, $\TMstroke$ के खंडों द्वारा, निरूपण होता है।) इसके बाद
  वह आगत ($n$ संख्या के~$\TMstroke$ का खंड) को दाईं ओर प्रतिलिपित करती
  है---किंतु प्रत्येक $\TMstroke$ के स्थान पर तीन $\TMstroke$ का एक खंड
  रखती है (याद रखें, $\TMstroke$ प्रतीक की संख्या~$3$ है; $1$,
  $\TMendtape$ की संख्या है और $2$, $\TMblank$ की संख्या है)। खंडों
  के इस अनुक्रम के बाएँ सिरे पर (जो $\TMblank$ प्रतीकों द्वारा~$U$ की
  पट्टी पर पृथक हैं) वह एक अकेला~$\TMstroke$ लिखती है, जो
  $\TMendtape$ का कूट है।

  अब~$U$ की पट्टी पर ये हैं: सूचकांक~$e$, संख्या~$n$, आरंभिक अवस्था की
  कूट-संख्या (``वर्तमान अवस्था''), आरंभिक शीर्ष-स्थान~$1$ की संख्या
  (``वर्तमान शीर्ष-स्थान''), और ``पट्टी'' की आरंभिक सामग्री
  ($\TMstroke$-खंडों का अनुक्रम, जो~$M_e$ के प्रतीकों की
  कूट-संख्याओं---``प्रतीकों''---को निरूपित करता है और~$\TMblank$ द्वारा
  पृथक है)।

  अब वह निम्न कार्यों द्वारा इसका अनुकरण करती है कि $M_e$ आगत~$n$ पर
  आरंभ होने पर क्या करती:
  \begin{enumerate}
    \item ``वर्तमान शीर्ष-स्थान'' की संख्या~$k$ खोजें (आरंभ में यह~$1$
    है),
    \item वहाँ कौन-सा ``प्रतीक'' है यह देखने के लिए ``पट्टी'' के $k$वें
    खंड पर जाएँ,
    \item\ollabel{find-inst}%
    वर्तमान ``अवस्था'' और ``प्रतीक'' से मेल खाने वाला निर्देश खोजें,
    \item ``पट्टी'' के $k$वें खंड पर वापस जाएँ और वहाँ के ``प्रतीक'' को
    उस प्रतीक की कूट-संख्या से बदलें जिसे $M_e$ लिखती,
    \item शीर्ष को उस स्थान पर ले जाएँ जहाँ वर्तमान ``अवस्था'' दर्ज है
    और वहाँ की संख्या को नई अवस्था की संख्या से बदलें,
    \item उस स्थान पर जाएँ जहाँ ``पट्टी-स्थान'' दर्ज है और एक
    ~$\TMstroke$ मिटाएँ अथवा एक~$\TMstroke$ जोड़ें (यदि निर्देश क्रमशः
    बाईं अथवा दाईं ओर जाने को कहता है)।
    \item दोहराएँ।\footnote{हम यहाँ कुछ सूक्ष्म कठिनाइयों का विस्तार
    छोड़ रहे हैं। उदाहरणतः ``वर्तमान शीर्ष-स्थान'' का लेखा रखने वाले
    गणक को बढ़ाते समय~$U$ को कुछ अतिरिक्त स्थान चाहिए हो सकता है---उस
    स्थिति में उसे पूरी ``पट्टी'' को दाईं ओर खिसकाना होगा।}
  \end{enumerate}
  यदि $M_e$ आगत~$n$ पर आरंभ होकर कभी नहीं रुकती, तो~$U$ भी कभी नहीं
  रुकती और इसलिए उसका निर्गत अपरिभाषित है।

  यदि चरण~\olref{find-inst} में पता चलता है कि~$M_e$ के वर्णन में
  वर्तमान ``अवस्था''/``प्रतीक'' युग्म के लिए कोई निर्देश नहीं है, तो
  $M_e$ रुक जाती। ऐसा होने पर~$U$ अपनी पट्टी का ``पट्टी'' के बाईं ओर का
  भाग मिटा देती है। तीन~$\TMstroke$ के प्रत्येक खंड के लिए (जो एक
  ~$\TMstroke$, $M_e$ की पट्टी पर, निरूपित करता है), वह अपनी पट्टी के बाएँ सिरे
  पर एक $\TMstroke$ लिखती है और क्रमशः ``पट्टी'' को मिटाती जाती है। यह
  पूरा होने पर~$U$ की पट्टी में~$\TMstroke$ का एक अकेला खंड होता है
  जिसकी लंबाई~$m$ है।
  
  यदि~$U$ को ``पट्टी'' पर तीन~$\TMstroke$ के खंड के अतिरिक्त कुछ और
  मिलता है, तो वह तुरंत रुक जाती है। चूँकि इस स्थिति में~$U$ की पट्टी
  में~$\TMstroke$ का एक अकेला खंड नहीं होता, उसका निर्गत कोई प्राकृतिक
  संख्या नहीं है; अर्थात् इस स्थिति में $f(e,n)$ अपरिभाषित है।
\end{proof}

\end{document}
